% Spoken scripts for the three main screencasts.  Included by screencast_guide.tex.
% \act{...}  what the presenter does on screen (not spoken)
% \say{...}  what the presenter says

\clearpage
\section*{Script: LessonsLearnedScreencast}
\scriptmeta{Target 6 minutes (allowed 5--10). Arush Ugal about 3.5 minutes, then Anirudh B about
2.5 minutes. Before recording: open \texttt{Lessons\_Learned\_Javadocs/index.html} in Chrome at 125\%,
open the folder in Finder, and open this guide at the Task 3 table.}

\speaker{Arush Ugal}{about 3.5 minutes}

\act{Show this guide, the ``LessonsLearnedScreencast: Task 3 items'' table.}
\say{This is the LessonsLearned screencast for Team 24. I'm Arush Ugal, and I'm presenting with
Anirudh B. We're showing the Lessons\_Learned\_Javadocs folder, which documents the two entity
classes, LessonsLearned and LessonsLearnedList. This table is the list of deliverable
requirements, 3.1 to 3.5, and next to each one are the HTML pages and screenshots that satisfy
it. We'll go through them in that order.}

\act{Open \texttt{3.1-entityClasses-package-source-tree.png}, then
\texttt{3.1-entityClasses-package-declarations.png}, then \texttt{3.1-entityClasses-package-summary-javadoc.png}.}
\say{Requirement 3.1 is that both classes are in the entityClasses package. The first screenshot
is the source tree: LessonsLearned.java and LessonsLearnedList.java sit in entityClasses, next
to the supplied User class. The second shows line 1 of each file, ``package entityClasses''. The
third is the generated package page, which lists both classes and describes what the package
is for.}

\act{Open \texttt{entityClasses/LessonsLearned.html}; scroll through the class description.}
\say{This is the Javadoc for LessonsLearned. One object is one lesson. The description says what
the class is for, why three attributes are stored now for TP3, and then a list that maps each
Contributor requirement, CS-1 to CS-12, to the method that implements it. That mapping is
requirement 3.4. The screenshot \texttt{3.4-LessonsLearned-class-javadoc.png} shows the same text in the
source.}

\act{Scroll to Field Summary; open \texttt{3.3-LessonsLearned-attributes-contributor-stories.png}
and \texttt{3.3-LessonsLearned-attributes-TP3-epics.png}.}
\say{Requirement 3.3 is the attributes. The Field Summary lists all twelve, including the
private ones, because the Javadoc was generated with the private option. In the source, each
attribute starts with ``Source'', the requirement it comes from, and then the reason it's stored
that way. For example, contributor comes from CS-1, CS-6, and CS-9: it's taken from the logged-in
user and never typed, so nobody can post as someone else. The second screenshot shows the three
TP3 attributes: review status, reviewer comment, and helpful count. We store them now so the
table doesn't have to change once it already holds lessons.}

\act{Scroll to Method Summary; open the three \texttt{3.2-LessonsLearned-*} screenshots in turn.}
\say{Requirement 3.2 is CRUD. Create validates first, refuses a lesson that's already stored,
sets both dates, and writes the row. Read is static, because there's no object until the row is
read. Update runs the same validation as create, and if the write fails it puts the old date
back, so the object never claims a change the database doesn't have. Delete sets the id back to
minus one. Validate is the single call that both create and update make.}

\act{Open \texttt{3.4-LessonsLearned-constructors.png}, then the two \texttt{3.5-LessonsLearned-*} screenshots.}
\say{There are two constructors: one for a new lesson, with no id and no dates, and one for a
row read back from the database. Requirement 3.5 is comments that explain code that isn't
obvious. This block explains why there's no setter for the contributor, the dates, or the TP3
attributes. And this one explains why isOwnedBy compares exactly, and why an empty search term
matches nothing. Anirudh will now cover LessonsLearnedList.}

\speaker{Anirudh B}{about 2.5 minutes}

\act{Open \texttt{entityClasses/LessonsLearnedList.html}, then \texttt{3.4-LessonsLearnedList-class-javadoc.png}
and \texttt{3.3-LessonsLearnedList-attribute.png}.}
\say{I'm Anirudh. LessonsLearnedList holds a list of lessons. readAll gives the persistent list
of every lesson in the system, and every other list is a subset of it. The class Javadoc explains
why a subset comes back as another LessonsLearnedList: you can narrow it again, count it, or show
it with the same code. It has one attribute, an ArrayList, declared final, with its source and
rationale written above it.}

\act{Open the two \texttt{3.2-LessonsLearnedList-*} screenshots and \texttt{3.2-LessonsLearnedDatabase-addLesson.png}.}
\say{For requirement 3.2, the list's CRUD is readAll, createLesson, updateLesson, and deleteLesson.
Each one writes to the database first and changes the in-memory list only if the write succeeded.
Underneath is LessonsLearnedDatabase. addLesson uses a PreparedStatement and reads back the key
the database generated.}

\act{Open \texttt{3.4-LessonsLearnedList-subsets.png} and \texttt{3.4-LessonsLearnedList-contributor-and-categories.png}.}
\say{These are the subsets: search, which covers title, description, and keywords and ignores
case; filter by category; filter by contributor, which uses the same isOwnedBy test as the
ownership rule; and getCategories, which fills the category filter on the list page.}

\act{Open the \texttt{3.5-LessonsLearnedDatabase-*}, \texttt{3.5-LessonInputPolicy-*}, and \texttt{3.5-LessonsLearnedList-getAll-returns-copy.png} screenshots.}
\say{The rest of requirement 3.5 is here. createTables explains why the new columns are added
with ALTER TABLE rather than in the CREATE statement. moveLegacyTableAside fixes a defect we
found: an HW2 table with different columns made every save fail with ``Column DESCRIPTION not
found''. Now that table is renamed and kept, never dropped. getAllLessons orders by id rather
than by date. updateLesson's comment explains which columns an edit can't change. checkLongText
records the HW2 defect that refused any description over 256 characters. And getAll returns a
copy, so a caller can't change the list behind its back. That's the Lessons\_Learned\_Javadocs
folder. Thanks for watching.}

% =====================================================================================
\clearpage
\section*{Script: PhaseTwoGUIScreencast}
\scriptmeta{Target 9 minutes (allowed 5--10). Karan Patel about 4.5 minutes (guiLessonsList and
the overall structure), then Moksha Dev about 4.5 minutes (guiLessonEditor and guiLessonDetail).
Before recording: open \texttt{Contributor\_Javadocs/index.html}, the folder in Finder, and this
guide at the Task 4 table. Have the application running for Moksha's last part.}

\speaker{Karan Patel}{about 4.5 minutes}

\act{Show this guide, the ``PhaseTwoGUIScreencast: Task 4 items'' table.}
\say{This is the Phase Two GUI screencast for Team 24. I'm Karan Patel, presenting with Moksha
Dev. We're showing the Contributor\_Javadocs folder. This table lists the five deliverable
requirements, 4.1 to 4.5, and next to each one the HTML pages and screenshot names that cover
it.}

\act{Open \texttt{4.1-Contributor-packages-overview-javadoc.png}, \texttt{4.1-gui-packages-source-tree.png},
then the three \texttt{4.1-*-package-classes-javadoc.png}.}
\say{For requirement 4.1 there are three new packages: guiLessonsList, guiLessonEditor, and
guiLessonDetail. They're named the same way as the supplied packages, like guiUserLogin and
guiRole1. The source tree shows each package has exactly three classes, a Model, a View, and a
Controller. The tree also marks the two home pages that got a Lessons Learned button. These
package pages list the classes in each one.}

\act{Open \texttt{4.2-guiLessonsList-View-wires-Controller.png}.}
\say{Requirement 4.2 is that we use MVC the way Foundations does. In the View's constructor,
every control gets a setOnAction that calls a Controller method: Search, Show All, the category
box, Only my lessons, New Lesson, Open, Back, and Quit. The View doesn't make any decisions.}

\act{Open \texttt{4.2-guiLessonsList-Controller-asks-Model.png}.}
\say{The Controller reads the controls, asks ModelLessonsList what to show and what the status
line should say, and puts the answer on screen. One difference from the supplied packages: there
the Model is empty, but ours holds data, because this page narrows a list and has to explain the
result. It has no JavaFX in it, so the tests can check its decisions directly.}

\act{Open \texttt{guiLessonsList/ModelLessonsList.html}, scroll to the Field Summary; then
\texttt{4.4-ModelLessonsList-attributes.png}, \texttt{4.4-ViewLessonsList-attributes.png}, and the three
\texttt{4.4-*LessonsList-class-javadoc.png}.}
\say{Requirement 4.4 is documenting purpose, attributes, and operations. ModelLessonsList has two
attributes, the full list and the subset on screen, each with its source requirement. The View's
attributes are the widgets, grouped by area. The class Javadoc for each of the three says what
the class is for, its role in MVC, and which requirements it supports: CS-2, CS-3, CS-4, CS-5,
and CS-6, plus the entry point to CS-1.}

\act{Open \texttt{4.5-ModelLessonsList-applyFilters.png}, \texttt{4.5-ModelLessonsList-statusMessage.png},
\texttt{4.5-ControllerLessonsList-reloadLessons-fillCategories.png}.}
\say{Requirement 4.5 is comments that explain code that isn't obvious. applyFilters explains why
an empty search box shows every lesson on this page, while the list class treats an empty term as
matching nothing. statusMessage explains why an empty result is described in words that name the
filter, so the user knows which filter to change. reloadLessons explains why the page re-reads the
database instead of keeping what's in memory: that's how lessons from an earlier run show up.}

\act{Open \texttt{4.4-guiRole1-Lessons-Learned-button.png} and the two \texttt{performLessonsLearned} screenshots.}
\say{The way in is the Lessons Learned button on the Role1 and Role2 home pages. Each one calls
performLessonsLearned, which passes the current stage and user along, so Back returns to the right
home page. Moksha will now cover the editor and the detail page.}

\speaker{Moksha Dev}{about 4.5 minutes}

\act{Open \texttt{guiLessonEditor/package-summary.html}, \texttt{4.4-guiLessonEditor-package-info.png},
\texttt{4.2-guiLessonEditor-View-and-Controller.png}.}
\say{I'm Moksha. guiLessonEditor is one page used for both a new lesson and an edit, because the
fields and rules are the same. The package description lists the requirements it covers: CS-1,
CS-7, CS-9, and the on-screen side of CS-10 and CS-11. For MVC, pressing Save Lesson calls
performSave in the Controller. That reads the fields, calls the Model, and if a field is refused,
shows it in a ValidationAlert.}

\act{Open \texttt{4.4-ViewLessonEditor-attributes.png}, \texttt{4.4-ModelLessonEditor-class-javadoc.png}.}
\say{Look at the View's attributes: the contributor is a Label, not a TextField, so it can't be
typed. Each field has a character counter for the CS-11 limits.}

\act{Open \texttt{4.5-ModelLessonEditor-saveNewLesson.png}, \texttt{4.5-ModelLessonEditor-saveChanges.png},
\texttt{4.5-ModelLessonEditor-tidyKeywords.png}.}
\say{For 4.5, saveNewLesson records why the contributor comes from the login: in HW2 it was a
text field, and that let anyone post in someone else's name. saveChanges checks ownership again,
even though the page only offers Edit to the contributor. It refuses an untouched form, and it
changes a copy rather than the original, so a refused save can't leave the list showing values
the database doesn't have. tidyKeywords drops duplicates but keeps an empty entry, so validation
can report it instead of the editor quietly changing what was typed.}

\act{Open \texttt{guiLessonDetail/package-summary.html}, \texttt{4.4-ViewLessonDetail-attributes.png},
\texttt{4.4-ViewLessonDetail-displayLessonDetail.png}, \texttt{4.4-ModelLessonDetail-canChange-ownershipNote.png}.}
\say{guiLessonDetail shows one lesson in full, which is CS-3, including the TP3 review line as
read-only. displayLessonDetail turns Edit and Delete on only for the contributor. The rule comes
from canChange in the Model, and ownershipNote explains on screen whose lesson it is. That's
CS-9.}

\act{Open \texttt{4.2-guiLessonDetail-View-and-Controller.png} and \texttt{4.5-ModelLessonDetail-deleteLesson.png}.}
\say{Delete Lesson calls performDelete, which asks for confirmation before it calls the Model.
That's CS-8. In deleteLesson, the comment explains why ownership is checked again in the Model:
the button is already disabled, but no other route to this method should be able to get around
the rule.}

\act{Open \texttt{4.5-FoundationsMain-lessonsDatabase.png}; then switch to the running application,
open the list, a lesson, and the editor for about 30 seconds.}
\say{Finally, FoundationsMain opens the lessons table when the application starts, and the
comment explains why a second connection to the same file doesn't trigger the ``already in use''
alert. Here are the three pages running: the list with search and filters, the detail page with
Edit and Delete, and the editor with its counters. That's the Contributor\_Javadocs folder.
Thanks for watching.}

% =====================================================================================
\clearpage
\section*{Script: TestingScreencast}
\scriptmeta{Target 12 minutes (allowed 8--15). Pranjal Ganvir about 6.5 minutes (the map, the
test documentation, and the Validation and CRUD runs), then Anirudh B about 5.5 minutes (List and
Model runs, totals, and the manual procedures). Before recording: open
\texttt{Requirements to Tests Map.pdf}, \texttt{ContributorTests\_Javadocs/index.html}, Eclipse with
the project, and the application ready to start with an empty database. In Eclipse, increase the
console font so the Pass lines can be read.}

\speaker{Pranjal Ganvir}{about 6.5 minutes}

\act{Show this guide, the ``TestingScreencast: Task 7 items'' table, then
\texttt{Requirements to Tests Map.pdf} pages 1 and 2.}
\say{This is the Testing screencast for Team 24. I'm Pranjal Ganvir, presenting with Anirudh B.
This table is the list of deliverable requirements, 7.1 to 7.4, with the files for each. For 7.1,
the Requirements to Tests Map lists all twelve requirements, CS-1 to CS-12. For 7.2, each row
gives its positive tests, its negative tests, and the HTML pages that document them. Every
requirement has at least one of each. There are eighty-three tests: seventy-six automated tests
in four suites, and seven manual procedures for things only a window can show.}

\act{Scroll to ``Every test, its Javadoc, and its screenshot''.}
\say{This table lists every test, its requirement, whether it's positive or negative, its method
in the HTML page, and the name of its screenshot.}

\act{Open \texttt{testingAutomation/LessonsCrudTestingAutomation.html}, Method Detail for
\texttt{crud11And12PersistAcrossReopen()}; then \texttt{7.4-CRUD-11-CRUD-12-PersistAcrossReopen.png}.}
\say{For 7.3, every test is its own method, and its Javadoc says three things: the requirement
and whether the test is positive or negative, how the test works, and what output means it
passed. This pair proves CS-12. It writes two lessons to a database file, deletes one, closes the
connection, which is what quitting the app does, then opens a new connection to the same file,
which is what restarting does. For 7.4, this is that method's screenshot, with the comments that
mark the first run and the second run.}

\act{Open \texttt{7.4-check-helper-and-refusedWith.png}.}
\say{Every test prints its result through check: ``Pass'' or ``Failure'', then the test id and the
requirement. For validation, a negative test only passes if the refusal message is the one the
requirement calls for, so a value refused for the wrong reason counts as a Failure.}

\act{In Eclipse, run \texttt{testingAutomation.TP2TestRunner}. Scroll to the Validation block.}
\say{Here's the live run. Validation, CS-10 and CS-11. VAL-01 to VAL-05 are the title: a normal
title passes; a blank one and one made only of spaces are refused as required; eighty characters
passes and eighty-one is refused as too long. VAL-06 to 09 do the same for category, at thirty
and thirty-one characters. VAL-10 to 14 are the description. VAL-12, three hundred characters,
is the regression test for the HW2 defect that refused anything over 256. VAL-13 accepts exactly
two thousand and VAL-14 refuses two thousand and one. VAL-15 to 17 show resolution is optional
and capped at two thousand. VAL-18 to 21 show keywords are optional, that four short ones pass,
that eleven are refused, and that a keyword over thirty characters is refused. VAL-22 and 23 are
the contributor. VAL-24 accepts a complete lesson, and VAL-25 shows a whole lesson is refused when
only its title is blank. Twenty-five passed, none failed.}

\act{Scroll to the CRUD block.}
\say{CRUD, which covers create, read, update, delete, ownership, and persistence. CRUD-01 stores a
valid lesson, gets an id, and the row count goes up by one. CRUD-02 refuses a blank title and the
count stays the same, so nothing was written. CRUD-03 reads that lesson back with all its fields,
and CRUD-04 shows an unknown id returns nothing instead of an error. CRUD-05 changes the title and
re-reads it from the database: the new title is there, and the contributor and creation date are
unchanged. CRUD-06 refuses an update to a blank title and the stored title stays put. CRUD-07
deletes the lesson, and CRUD-08 shows deleting it again reports false. CRUD-09 and 10 check
ownership: the contributor is recognised; another user, the same name in different capitals, and
a missing name are not. CRUD-11 and 12 are the restart test I just showed. CRUD-13 and 14 build
the table layout HW2 left behind: saving now works, and the old HW2 row is kept in the renamed
table. Fourteen of fourteen. Anirudh will take the List and Model suites.}

\speaker{Anirudh B}{about 5.5 minutes}

\act{Scroll to the List block.}
\say{I'm Anirudh. The List suite works on three seeded lessons: A and B in Database, C in GUI.
LIST-01 reads all three, and LIST-02 shows an empty database gives an empty list rather than an
error. LIST-03 searches ``JDBC'', which only appears in keywords, and finds exactly A and B, so
keywords are searched. LIST-04 searches a word no lesson has and gets nothing. LIST-05 filters
the Database category, typed in lower case, and gets A and B. LIST-06 gets nothing for a category
nobody uses. LIST-07 returns exactly Arush's two lessons, and LIST-08 nothing for an unknown
user. LIST-09 gives the category filter its choices, Database and GUI, each once. LIST-10 chains a
search and a filter. LIST-11 confirms newest first, C, B, A. LIST-12 shows search ignores case.
LIST-13 shows an empty search term matches nothing at the list level. Thirteen of thirteen.}

\act{Scroll to the Model block.}
\say{The Model suite tests the decisions the three windows make. MOD-01 saves through the editor
and checks the contributor is the logged-in user. MOD-02 refuses a blank title. MOD-03 shows that
with no filters the page shows everything: ``Showing 2 of 2 lessons''. MOD-04 checks a user with no
lessons is told so in words. MOD-05 and 06 are search on the page, and MOD-17 checks the no-match
message names the search term. MOD-07 and 08 are the category filter, and MOD-18 checks the
message names the category. MOD-09 and 10 are ``Only my lessons''. MOD-11 and 12 are the
ownership rule behind the Edit and Delete buttons. MOD-13 saves the contributor's change; MOD-14
refuses an edit that empties the description. MOD-15 deletes, and the lesson is gone from the
database and the shared list. MOD-16 refuses a delete by someone else, and MOD-20 refuses their
edit, with the stored title unchanged. MOD-19 is the message for an empty database. MOD-21 to 24
came from Moksha's review: a keyword list starting or ending with a comma is refused instead of
silently changed, ``All categories'' can't be used as a category name because the filter uses it,
and names that only look similar, like ``All'', still work. Twenty-four of twenty-four.}

\act{Scroll to the combined block.}
\say{The combined result: seventy-six tests, seventy-six passed, none failed. To make sure these
tests can actually fail, we put three defects into a copy of the code, and six tests failed, each
one naming the requirement that broke.}

\act{Start \texttt{LessonsManualTestProcedures} beside the running application (empty database,
users contribA and contribB already created). Carry out each procedure and type P.}
\say{The manual procedures cover what only a window can show. MAN-01: as contribA, the empty list
says no lessons have been recorded yet. MAN-02: in New Lesson the contributor is a label, not a
field, and the saved lesson shows ``by contribA''. MAN-03: a blank title brings up ``The title is
required'' and nothing is saved. MAN-04: Delete asks first; Cancel keeps the lesson and OK
removes it. MAN-05: as contribB, Edit and Delete are disabled and the note names contribA.
MAN-06: contribB can't find the deleted lesson. MAN-07: after quitting and restarting, the
lesson is still there. Seven passed. That's the ContributorTests\_Javadocs folder. Thanks for
watching.}
